% 日本語サンプル（lualatex + luatexja が必要。Debian/Ubuntu では texlive-lang-japanese）
\documentclass[11pt]{ltjsarticle}
\usepackage{amsmath}
\usepackage[margin=25mm]{geometry}

\title{日本語 LaTeX のプレビュー確認}
\author{html-preview-hub サンプル}
\date{}

\begin{document}
\maketitle

\section{目的}
AI に生成させた日本語の \LaTeX{} 文書が、アプリ内でそのまま PDF として読めることを確認する。
エンジンは \verb|\documentclass{ltjsarticle}| から lualatex と推定される。

\section{数式}
ソフトマックスの温度 $\tau$ を含む形は次のとおり。
\begin{equation}
  p_i = \frac{\exp(z_i / \tau)}{\sum_j \exp(z_j / \tau)}
\end{equation}

\section{表}
\begin{center}
\begin{tabular}{lr}
  \hline
  項目 & 値 \\
  \hline
  スキャン時間 & 134 ms \\
  初期表示 & 188 ms \\
  \hline
\end{tabular}
\end{center}

\end{document}
